\chapter{常用符号与记号表}
\label{app:notation}

本附录汇总本书与后续模形式学习中反复出现的常用符号，便于查阅。表中“首次出现章节”按本书结构标注：凡属于线性代数、群论、环论、模论、域论与伽罗瓦理论的记号，均对应本书主文第~1--12~章；凡直接服务于模形式正文而尚未在本书主文中展开定义的记号，统一标为“后续模形式章节”；范畴论相关记号则见附录~A。

\begingroup
\setlength{\LTleft}{0pt}
\setlength{\LTright}{0pt}
\renewcommand{\arraystretch}{1.2}

\begin{longtable}{@{}>{\raggedright\arraybackslash}p{0.17\textwidth}>{\raggedright\arraybackslash}p{0.18\textwidth}>{\raggedright\arraybackslash}p{0.41\textwidth}>{\raggedright\arraybackslash}p{0.14\textwidth}@{}}
\caption{常用符号与记号总表}\\
\toprule
符号 & 中文名称 & 简要说明 & 首次出现章节 \\
\midrule
\endfirsthead

\multicolumn{4}{c}{\tablename\ \thetable\ （续）}\\
\toprule
符号 & 中文名称 & 简要说明 & 首次出现章节 \\
\midrule
\endhead

\midrule
\multicolumn{4}{r}{续下页}\\
\endfoot

\bottomrule
\endlastfoot

\multicolumn{4}{@{}l}{\textbf{一、数域与集合}}\\
\addlinespace[0.3em]
$\mathbb{N}$ & 自然数集 & 默认表示正整数集合；若上下文需要包含 $0$，正文会另作说明。 & 第1章 \\
$\mathbb{Z}$ & 整数集 & 全体整数所成的交换环，也是最基本的整环例子。 & 第1章 \\
$\mathbb{Q}$ & 有理数域 & 由整数分式构成的最小特征零数域。 & 第1章 \\
$\mathbb{R}$ & 实数域 & 线性代数与分析背景中最常用的基域之一。 & 第1章 \\
$\mathbb{C}$ & 复数域 & 代数闭域；在特征值理论、内积空间与模形式中都极为重要。 & 第1章 \\
$\mathbb{F}_q$ & $q$ 元有限域 & 含有 $q=p^r$ 个元素的有限域，常用于有限域、表示与伽罗瓦理论。 & 第11章 \\
$\mathbb{H}$ & 上半平面 & 满足 $\operatorname{Im}(z)>0$ 的复数全体，是经典模形式的定义域。 & 后续模形式章节 \\
$\varnothing$ & 空集 & 不含任何元素的集合；用于子集、交并运算与拓扑论述。 & 第1章 \\
\addlinespace[0.6em]

\multicolumn{4}{@{}l}{\textbf{二、矩阵与线性代数}}\\
\addlinespace[0.3em]
$\operatorname{Mat}_{m\times n}(F)$ & $F$ 上的矩阵空间 & 所有 $m\times n$ 的 $F$ 值矩阵所成的集合，通常也视作向量空间。 & 第1章 \\
$\mathrm{GL}_n(F)$ & 一般线性群 & 由所有可逆的 $n$ 阶 $F$ 矩阵组成的群，群运算为矩阵乘法。 & 第7章 \\
$\mathrm{SL}_n(F)$ & 特殊线性群 & $\mathrm{GL}_n(F)$ 中行列式等于 $1$ 的子群。 & 第7章 \\
$O(n)$ & 正交群 & 保持实欧氏内积不变的线性变换群。 & 第7章 \\
$\mathrm{SO}(n)$ & 特殊正交群 & 正交群中行列式为 $1$ 的连通部分，在几何与表示论中常见。 & 第7章 \\
$U(n)$ & 酉群 & 保持厄米内积不变的复线性变换群。 & 第7章 \\
$\mathrm{SU}(n)$ & 特殊酉群 & 酉群中行列式为 $1$ 的子群。 & 第7章 \\
$\det(A)$ & 行列式 & 刻画线性变换体积伸缩与可逆性的标量不变量。 & 第1章 \\
$\operatorname{tr}(A)$ & 迹 & 方阵对角元之和；与特征值之和、表示角色密切相关。 & 第4章 \\
$\operatorname{rank}(A)$ & 秩 & 矩阵列空间或行空间的维数，衡量线性方程组的独立信息量。 & 第3章 \\
$\operatorname{diag}(a_1,\dots,a_n)$ & 对角矩阵记号 & 以 $a_1,\dots,a_n$ 为主对角元、其余元素为 $0$ 的矩阵。 & 第1章 \\
$I_n$ & $n$ 阶单位矩阵 & 主对角线上全为 $1$、其余元素全为 $0$ 的方阵。 & 第1章 \\
$A^T$ & 转置矩阵 & 由交换矩阵的行与列得到的矩阵。 & 第1章 \\
$A^*$ & 共轭转置或伴随 & 在线性内积空间中表示与内积相容的伴随算子矩阵。 & 第5章 \\
$A^{-1}$ & 逆矩阵 & 满足 $AA^{-1}=A^{-1}A=I_n$ 的矩阵，仅在 $A$ 可逆时存在。 & 第1章 \\
$\operatorname{adj}(A)$ & 伴随矩阵 & 由余子式矩阵转置得到，用于公式 $A^{-1}=\det(A)^{-1}\operatorname{adj}(A)$。 & 第1章 \\
\addlinespace[0.6em]

\multicolumn{4}{@{}l}{\textbf{三、群论}}\\
\addlinespace[0.3em]
$|G|$ & 群的阶 & 若 $G$ 有限，则表示其元素个数；更一般时也可表示基数。 & 第6章 \\
$[G:H]$ & 指数 & 子群 $H\leq G$ 的左陪集个数；有限群中满足 $|G|=[G:H]|H|$。 & 第6章 \\
$G/N$ & 商群 & 当 $N\trianglelefteq G$ 时，由 $N$ 的陪集组成的群。 & 第6章 \\
$Z(G)$ & 群的中心 & 与 $G$ 中每个元素都可交换的元素所成子群。 & 第6章 \\
$\operatorname{Aut}(G)$ & 自同构群 & 群 $G$ 到自身的全部群同构所成之群。 & 第6章 \\
$\operatorname{Inn}(G)$ & 内自同构群 & 由共轭作用 $x\mapsto gxg^{-1}$ 产生的自同构组成的正规子群。 & 第6章 \\
$\operatorname{Syl}_p(G)$ & 西罗 $p$ 子群族 & 有限群中阶为 $p^r$ 的极大 $p$ 子群的全体或其任一代表。 & 第6章 \\
$S_n$ & 对称群 & $n$ 个元素全排列组成的群。 & 第6章 \\
$A_n$ & 交错群 & 对称群中全部偶置换组成的正规子群。 & 第6章 \\
$D_n$ & 二面体群 & 正 $n$ 边形的全体对称变换组成的群，常取其阶为 $2n$。 & 第6章 \\
$C_n$ & 循环群 & 由一个元素生成、阶为 $n$ 的阿贝尔群。 & 第6章 \\
\addlinespace[0.6em]

\multicolumn{4}{@{}l}{\textbf{四、环与模}}\\
\addlinespace[0.3em]
$R[x]$ & 多项式环 & 以交换环 $R$ 为系数环的一元多项式环。 & 第9章 \\
$R/I$ & 商环 & 由环 $R$ 对理想 $I$ 取商得到的环。 & 第9章 \\
$(a)$ & 主理想 & 由单个元素 $a$ 生成的理想，即 $aR$ 或 $Ra$。 & 第9章 \\
$\operatorname{Ann}(M)$ & 湮没子 & 使模 $M$ 中每个元素都变为零的环元素所成理想。 & 第10章 \\
$\operatorname{Rad}(I)$ & 理想根 & 所有某次幂落入 $I$ 的元素组成的理想。 & 第9章 \\
$\operatorname{char}(R)$ & 环的特征 & 使 $n\cdot 1_R=0$ 的最小正整数；若不存在则记为 $0$。 & 第9章 \\
$\operatorname{Spec}(R)$ & 素谱 & 环 $R$ 的全部素理想所成的集合，也是交换代数几何的基本对象。 & 第9章 \\
$M\otimes_R N$ & 张量积 & 两个 $R$ 模的双线性信息所表示的泛对象。 & 第10章 \\
$\operatorname{Hom}_R(M,N)$ & 模同态集 & 从 $M$ 到 $N$ 的所有 $R$ 线性映射组成的集合或模。 & 第10章 \\
$\operatorname{End}_R(M)$ & 自同态环 & 模 $M$ 到自身的全部 $R$ 线性自映射所成之环。 & 第10章 \\
\addlinespace[0.6em]

\multicolumn{4}{@{}l}{\textbf{五、域论与伽罗瓦理论}}\\
\addlinespace[0.3em]
$[L:K]$ & 扩张次数 & 把域扩张 $L/K$ 视作 $K$ 向量空间时的维数。 & 第11章 \\
$\operatorname{Gal}(L/K)$ & 伽罗瓦群 & 固定基域 $K$ 的 $L$ 上域自同构全体组成的群。 & 第12章 \\
$\mathrm{Frob}_{\mathfrak p}$ & 弗罗贝尼乌斯元 & 在未分歧素理想 $\mathfrak p$ 处刻画有限域提升的特殊自同构。 & 第12章 \\
$\Phi_n(x)$ & 第 $n$ 个分圆多项式 & 以本原 $n$ 次单位根为根的既约多项式。 & 第12章 \\
$\mathbb{F}_q^\times$ & 有限域乘法群 & 有限域中全部非零元素组成的循环群。 & 第11章 \\
\addlinespace[0.6em]

\multicolumn{4}{@{}l}{\textbf{六、模形式}}\\
\addlinespace[0.3em]
$M_k(\Gamma)$ & 权 $k$ 模形式空间 & 关于群 $\Gamma$ 的权 $k$ 模形式组成的复向量空间。 & 后续模形式章节 \\
$S_k(\Gamma)$ & 权 $k$ 尖点形式空间 & 在各尖点处消失的权 $k$ 模形式所成子空间。 & 后续模形式章节 \\
$\Gamma_0(N)$ & $\Gamma_0(N)$ 同余子群 & 模 $N$ 意义下左下角元素可被 $N$ 整除的经典同余子群。 & 后续模形式章节 \\
$\Gamma_1(N)$ & $\Gamma_1(N)$ 同余子群 & 比 $\Gamma_0(N)$ 更细的同余子群，对角元模 $N$ 同余于 $1$。 & 后续模形式章节 \\
$\Gamma(N)$ & 主同余子群 & 从 $\mathrm{SL}_2(\mathbb{Z})$ 到 $\mathrm{SL}_2(\mathbb{Z}/N\mathbb{Z})$ 的核。 & 后续模形式章节 \\
$T_n$ & 赫克算子 & 作用在模形式空间上的一族两两可交换线性算子。 & 后续模形式章节 \\
$\langle f,g\rangle$ & 彼得松内积 & 模形式空间上的标准内积，用于正交分解与谱理论。 & 后续模形式章节 \\
$E_k$ & 艾森斯坦级数 & 最基本的一类模形式，常作为模形式空间的显式基元。 & 后续模形式章节 \\
\addlinespace[0.6em]

\multicolumn{4}{@{}l}{\textbf{七、范畴论}}\\
\addlinespace[0.3em]
$\operatorname{Ob}(\mathcal{C})$ & 对象类 & 范畴 $\mathcal{C}$ 中全部对象的集合或类。 & 附录A \\
$\operatorname{Mor}(\mathcal{C})$ & 态射类 & 范畴 $\mathcal{C}$ 中全部态射的集合或类。 & 附录A \\
$\operatorname{Hom}_{\mathcal{C}}(X,Y)$ & 态射集 & 从对象 $X$ 到对象 $Y$ 的全部态射组成的集合。 & 附录A \\
$F\colon \mathcal{C}\to \mathcal{D}$ & 函子 & 在两个范畴之间同时保持对象、态射及复合结构的映射。 & 附录A \\
\addlinespace[0.6em]

\multicolumn{4}{@{}l}{\textbf{八、逻辑与通用}}\\
\addlinespace[0.3em]
$\forall$ & 任意、对一切 & 全称量词，表示对讨论范围内每个对象都成立。 & 第1章 \\
$\exists$ & 存在 & 存在量词，表示至少存在一个对象使断言成立。 & 第1章 \\
$\Longrightarrow$ & 推出 & 表示由前一条件可推出后一结论。 & 第1章 \\
$\Longleftrightarrow$ & 当且仅当 & 表示两个命题彼此推出、逻辑上等价。 & 第1章 \\
$:=$ & 定义为 & 用于引入新记号、新对象或新运算的定义。 & 第1章 \\
$\cong$ & 同构 & 表示两个代数对象在相应范畴中结构完全相同。 & 第1章 \\
$\simeq$ & 等价或自然同构 & 表示依上下文而定的等价关系、弱同构关系或自然识别。 & 附录A \\
$\trianglelefteq$ & 正规子群符号 & 写作 $N\trianglelefteq G$，表示 $N$ 是 $G$ 的正规子群。 & 第6章 \\
$\oplus$ & 直和 & 用于向量空间、阿贝尔群、模或表示的直和分解。 & 第2章 \\
$\otimes$ & 张量积符号 & 既可表示对象层面的张量积，也可表示纯张量。 & 第10章 \\
$\prod$ & 乘积 & 可表示数乘、连乘、直积或范畴中的积对象。 & 第1章 \\
$\coprod$ & 余积 & 在集合论中常对应不交并，在范畴论中表示余积对象。 & 附录A \\
\end{longtable}
\endgroup
